\subsection{Three arms on a thousand items}\label{subsec:threearm}
The largest test compares three arms and nothing else. In the first the agent has the brain: the recall tool, served without a shell, and the two prompt hooks of Section~\ref{sec:hook}, and no file tools. In the second it has no brain, no hooks and no recall tool, but it can read, list and search the research repositories, which include the exported store, the campaign records and the open-problem files; the brain's own folders, its transcripts, the keys and every test folder are denied to it. In the third it has the model alone, with no tools. Every arm sees the same system text, the same instructions on how to answer and the same item, one item a session. Claude Sonnet agents (high reasoning) had $600$ seconds an item and Codex agents (gpt-6.1-sol, high reasoning) $900$ seconds. No Claude arm had a shell. For Codex the brain arm had the recall tool and the hook notes and no shell, the documents arm a read-only shell over the same repositories, and the third arm nothing. The runs took place on 23 and 24 October 2026. The Claude run splits the bank into three shards run side by side, and the Codex run takes items the Claude run has already finished; the figures below cover the items that had been run in all three arms by the evening of 24 October, $755$ of the $1{,}000$ for Claude and $182$ for Codex.

The bank was built on 12 September 2026 and has $1{,}000$ items. $923$ are lookups over the store's records: questions on the open-problem ledger (the sum or the difference of the pages on which listed problems are posed in their sources, which of four is posed latest, which source is the odd one out) and on the self-avoiding-walk counts held in the store, each combining two to four records. The other $77$ are proof audits, in which the agent decides whether a numbered proof is correct as written and, if not, names its first invalid step. The lookups test what the memory is for; the audits are a mathematics control, since their answers are not stored anywhere. A timeout or an error without an answer counts as wrong. A session that never reached the model (a refused account or a session limit, recognised by an empty event stream) is dropped and rerun, not graded.

\begin{table}[ht]
\centering\small
\begin{tabular}{llccc}
\toprule
 & & Brain & Documents & No tools\\
\midrule
Claude, $755$ items run in all arms & correct & $94.3\%$ & $78.1\%$ & $21.7\%$\\
 & sessions timed out & $4.8\%$ & $21.2\%$ & $35.8\%$\\
 & mean seconds an item & $136$ & $241$ & $249$\\
Claude, $593$ items both brain and documents & correct & $99.2\%$ & $99.3\%$ & \\
\quad answered in time & & & & \\
Claude, $480$ items all arms answered in time & correct & $99.2\%$ & $99.4\%$ & $34.0\%$\\
\quad of which proof audits, $68$ & correct & $98.5\%$ & $97.1\%$ & $98.5\%$\\
Codex, $182$ items run in all arms & correct & $93.4\%$ & $97.3\%$ & $31.9\%$\\
\quad of which proof audits, $50$ & correct & $98.0\%$ & $94.0\%$ & $92.0\%$\\
\bottomrule
\end{tabular}
\caption{The three-arm test, Claude Sonnet and Codex gpt-6.1-sol, 24 October 2026. ``Documents'' is the arm with the repositories and no brain.}
\label{tab:threearm}
\end{table}

For Claude the brain beat the documents $134$ to $12$ on the items where they disagreed ($p=3.0\cdot10^{-27}$) and the bare model $549$ to $1$; the documents beat the bare model $427$ to $1$. The first of these numbers needs reading with the rest of the table. On the $593$ items that the brain arm and the documents arm both answered within the time, they are tied, $99.2\%$ against $99.3\%$, $2$ to $3$. The brain's lead over the documents is therefore almost entirely time: an agent searching the files finds the same records, but it needs more searches to find them, and on about a fifth of the items it is still searching when the clock runs out. The gap is largest on the walk lookups that need three or four records, where the documents arm answered $31$ to $51\%$ of the items in a family and the brain arm $76$ to $95\%$. All but one of the bare model's timed-out sessions fall on these walk lookups, where it tries to enumerate the walks by hand.

The Codex agents had half again as much time and seldom ran out of it. There the documents arm scored $97.3\%$ against $93.4\%$ for the brain ($3$ to $10$, $p=0.09$; on the lookups alone $1$ to $10$, $p=0.012$, which does not survive the correction of Appendix~\ref{app:stats}). On the first $100$ Codex items every one of the brain arm's losses was an abstention, never a wrong value. In two the agent never called the recall tool. In three its searches missed records that the documents arm found by reading the files; in one of these the record held a string of punctuation and digits, which the recall tool's lexical match handles poorly and a text search finds at once. Neither aid moved the proof audits, which every arm, the bare model included, answered at about the same rate: the store adds nothing to a proof that the model can check by itself.

\subsubsection*{What it cost}
The $923$ lookups were generated by scripts from the store's records, and every key was derived twice, from separate reads, with no disagreement. The $77$ proof audits were written by Claude Opus~5 agents on 12 September 2026 and each key was confirmed by a separate verifier; by the workflow's own record that took $2.3$ million tokens and $32$ minutes. The record keeps the total but not its split between input and output, so the price can only be bounded: at list price, \$5 a million input tokens and \$25 a million output tokens, it lies between \$12 and \$58, and cached input would bring it lower. Table~\ref{tab:threearm-cost} gives the cost of every session of the three-arm runs, timeouts and reruns included. The prices are those the Claude command-line client reports at list price; the runs used subscription accounts, and the Codex client reports no price. A session stopped at its time limit records neither tokens nor a price, so the token and price columns leave out the timed-out sessions.

\begin{table}[ht]
\centering\small
\begin{tabular}{llrrrrr}
\toprule
 & Arm & Sessions & Hours & \multicolumn{2}{c}{Tokens a session} & Price a\\
\cmidrule(lr){5-6}
 & & & & in & out & session\\
\midrule
Claude & brain & $808$ & $30.6$ & $20{,}000$ & $1{,}500$ & \$0.053\\
 & documents & $808$ & $53.2$ & $80{,}000$ & $4{,}000$ & \$0.092\\
 & no tools & $807$ & $54.8$ & $2{,}000$ & $300$ & \$0.010\\
Codex & brain & $227$ & $27.6$ & $45{,}000$ & $400$ & ---\\
 & documents & $231$ & $17.0$ & $57{,}000$ & $600$ & ---\\
 & no tools & $232$ & $14.8$ & $7{,}000$ & $3{,}100$ & ---\\
\bottomrule
\end{tabular}
\caption{Cost of the three-arm runs to 24 October 2026: $3{,}113$ sessions, $198$ session-hours, $107$ million input tokens, $5.7$ million output tokens, and \$125.69 at list price for the Claude sessions.}
\label{tab:threearm-cost}
\end{table}

With Claude the brain arm read a quarter of the documents arm's input, wrote two fifths of its output and used $58\%$ of its session-hours. Counting only the sessions that recorded a price, it cost $58\%$ as much a session; the documents arm's timed-out sessions, four times as many as the brain arm's, are left out of that figure, so the true ratio is probably lower. With Codex the order of the hours is reversed: on the graded sessions the brain arm took $244$ seconds an item against $215$, and its total includes sessions rerun after the connection to the recall server dropped.
